\documentclass[11pt,a4paper]{article}
\usepackage[utf8]{inputenc}
\usepackage[T1]{fontenc}
\usepackage{geometry}
\geometry{top=1in, bottom=1in, left=1in, right=1in}
\usepackage{hyperref}
\usepackage{booktabs}
\usepackage{amsmath}
\usepackage{listings}
\usepackage{xcolor}
\usepackage{graphicx}
\usepackage{array}
\usepackage{fancyhdr}
\usepackage{titlesec}
\usepackage{longtable}

\hypersetup{
    colorlinks=true,
    linkcolor=purple,
    filecolor=magenta,      
    urlcolor=cyan,
    pdftitle=Genesis Observatory Specification (GOS),
    pdfpagemode=FullScreen,
}

\definecolor{codegray}{rgb}{0.96,0.96,0.96}
\definecolor{codeborder}{rgb}{0.85,0.85,0.85}
\definecolor{commentgreen}{rgb}{0.12,0.55,0.22}
\definecolor{keywordblue}{rgb}{0.1,0.2,0.8}

\lstset{
    backgroundcolor=\color{codegray},
    basicstyle=\ttfamily\small,
    breakatwhitespace=false,
    breaklines=true,
    captionpos=b,
    commentstyle=\color{commentgreen},
    frame=single,
    rulecolor=\color{codeborder},
    keepspaces=true,
    keywordstyle=\color{keywordblue},
    language=Python,
    showspaces=false,
    showstringspaces=false,
    showtabs=false,
    tabsize=4
}

\pagestyle{fancy}
\fancyhf{}
\rhead{\small Genesis Observatory Specification (GOS)}
\lhead{\small Project Genesis}
\cfoot{\thepage}

\title{\textbf{Genesis Observatory Specification (GOS)}}
\author{Project Genesis Research Repository}
\date{\small Generated: July 2026}

\begin{document}

\maketitle
\thispagestyle{empty}

\newpage
\section*{Genesis Observatory Specification (GOS)}
\subsection*{Version 1.0  July 2026}

The Observatory (Evidence Explorer) is the scientific instrument deck of the Genesis Portal. Its role is to explain the narrative: if the Civilization Record answers \textbackslash{}textit{what} happened, the Observatory answers \textbackslash{}textit{why} it happened using quantitative evidence.

\noindent\rule{\textwidth}{0.4pt}

\subsection*{1. Core Philosophy}

In accordance with \textbackslash{}href{file:///c:/Users/vinay/Desktop/project\%20genesis/docs/architecture/ADR/ADR-006-narrative-first.md}{ADR-006} and \textbackslash{}href{file:///c:/Users/vinay/Desktop/project\%20genesis/docs/architecture/ADR/ADR-007-goal-oriented-visualizations.md}{ADR-007}, the Observatory operates under two absolute constraints:
\begin{enumerate}
  \item \textbackslash{}textbf{Evidence supports narrative}: The Observatory is entered \textbackslash{}textit{from} the Civilization Record. It is a secondary explorer, not a standalone home page.
  \item \textbackslash{}textbf{Visualizations are Laboratory Instruments}: Every chart must answer a specific research question. Generic data plots without intent are prohibited.
\end{enumerate}

\noindent\rule{\textwidth}{0.4pt}

\subsection*{2. Instrument Catalog \& Registry}

We document and register all approved scientific instruments. Every instrument must map to a research question, inputs, and limitations.

\subsubsection*{OBS-001: Population Dynamics (Oscilloscope)}
\begin{itemize}
  \item \textbackslash{}textbf{Research Question}: \textbackslash{}textit{Did resource scarcity reduce survival and growth rates?}
  \item \textbackslash{}textbf{Inputs}: \textbackslash{}texttt{population.csv} (ticks, total count, and colony breakdowns).
  \item \textbackslash{}textbf{Outputs}: Line graph mapping total population and individual colony lineages over time.
  \item \textbackslash{}textbf{Limitations}: Does not show spatial distribution or local water pressures.
\end{itemize}

\subsubsection*{OBS-002: Resource Consumption Explorer (Flow Meter)}
\begin{itemize}
  \item \textbackslash{}textbf{Research Question}: \textbackslash{}textit{Where were critical resource shortages created?}
  \item \textbackslash{}textbf{Inputs}: \textbackslash{}texttt{summary.json} (derived metrics for food and water efficiency).
  \item \textbackslash{}textbf{Outputs}: Stacked area chart showing resource supply vs. total agent consumption.
  \item \textbackslash{}textbf{Limitations}: Aggregated globally; does not show individual agent hunger levels.
\end{itemize}

\subsubsection*{OBS-003: Territorial Friction Map (Microscope)}
\begin{itemize}
  \item \textbackslash{}textbf{Research Question}: \textbackslash{}textit{Did territorial pressure lead to clustered disputes?}
  \item \textbackslash{}textbf{Inputs}: \textbackslash{}texttt{events.json} (Dispute events metadata mapping coordinates).
  \item \textbackslash{}textbf{Outputs}: Static 2D spatial heatmap overlay showing conflict coordinates on the map.
  \item \textbackslash{}textbf{Limitations}: Ticks are collapsed; does not show the sequence of dispute progression.
\end{itemize}

\subsubsection*{OBS-004: Social Kinship Grid (Oscilloscope)}
\begin{itemize}
  \item \textbackslash{}textbf{Research Question}: \textbackslash{}textit{Did stable cooperative attachment bonds emerge?}
  \item \textbackslash{}textbf{Inputs}: \textbackslash{}texttt{agent\textbackslash{}\_census.csv} (parents, colonies, and children lineages).
  \item \textbackslash{}textbf{Outputs}: Dynamic directed network graph showing colony splits and child-parent attachments.
  \item \textbackslash{}textbf{Limitations}: Becomes visually complex when agent count exceeds 500.
\end{itemize}

\subsubsection*{OBS-005: Genetic Entropy Browser (Sequencer)}
\begin{itemize}
  \item \textbackslash{}textbf{Research Question}: \textbackslash{}textit{Did genetic traits stabilize or drift over generations?}
  \item \textbackslash{}textbf{Inputs}: \textbackslash{}texttt{agent\textbackslash{}\_census.csv} (chromosome strings and generation indices).
  \item \textbackslash{}textbf{Outputs}: Histogram showing distribution of mutation deviations from the founder seed.
  \item \textbackslash{}textbf{Limitations}: Captures survivors only; deceased lineages are excluded.
\end{itemize}

\noindent\rule{\textwidth}{0.4pt}

\subsection*{3. Interaction \& Linking Rules}

\begin{itemize}
  \item \textbackslash{}textbf{Tick Linking (Phase 3.5 Bridge)}: Clicking on any event in the Chronicle seeks all active instruments to that specific tick (e.g. updating the pointer on the Population Dynamics graph).
  \item \textbackslash{}textbf{Decoupled Tabs}: The Observatory loads in a split-screen or separate tab context, keeping the narrative chronicle layout readable.
\end{itemize}

\end{document}
